% U3 — KROS 총설 1부 (파생 원고; 정본은 ../../main.tex)
% 템플릿: A1 편집위 회신 후 KROS 공식 템플릿으로 이식 예정. 그 전까지 article.
\documentclass[11pt,a4paper]{article}
\usepackage{kotex}
\usepackage{amsmath,amssymb,bm}
\usepackage[margin=25mm]{geometry}
\usepackage{booktabs}
\usepackage{hyperref}
\newcommand{\dd}{\mathrm{d}}
\newcommand{\jj}{\mathrm{j}}

\title{임피던스의 역사와 기초: 전신선로에서 로봇 접촉 제어까지 --- 초보자를 위한 총설 (1부)}
\author{(저자 표기: USER\_ACTIONS C4 확정 후 기입)}
\date{}

\begin{document}
\maketitle

\begin{abstract}
협동로봇과 휴머노이드의 보급으로, 제어 전공자가 아닌 사용자도 로봇의 접촉
거동을 설정해야 하는 시대가 오고 있다.
접촉 거동의 핵심 언어인 임피던스는 그러나 대부분의 문헌에서 유도 과정이
생략된 채 전문가 눈높이로 서술되어, 입문자는 하나의 개념을 이해하기 위해
여러 문헌을 오가야 한다.
본 총설 1부는 임피던스라는 언어가 자란 역사적 문제(장거리 전신선로의 신호
왜곡)에서 출발하여, 라플라스 변환의 정의와 미분 규칙, RLC 회로의 전기
임피던스, 질량-스프링-댐퍼의 기계 임피던스까지를 유도 생략 없이 하나의
표기 체계로 정리한다.
모든 수치 예제는 공개 MuJoCo 교육 시뮬레이터에서 재현 가능하며, 세부 유도
전체를 담은 확장판 원고와 시뮬레이터는 공개 저장소로 제공한다.
2부에서는 이 기초 위에서 로봇 임피던스 제어와 파라미터 설정 원리를 다룬다.
\end{abstract}

\section{서론}
\label{sec:u3-intro}

산업 현장 밖으로 나온 로봇이 사람의 공간에서 물건을 집고, 닦고, 끼우는
일이 흔해지고 있다 \cite{nvidia_physical_ai_2025}.
이런 접촉 작업의 품질과 안전은 위치 정확도만으로 결정되지 않고, 로봇이
힘과 움직임 사이의 관계 --- 즉 임피던스 --- 를 어떻게 갖도록 설정되었는가에
좌우된다 \cite{hogan1985impedancepart1,abudakka2020variableimpedance}.
문제는 이 설정이 전공자에게도 쉽지 않다는 점이다.
강성과 목표 오프셋을 잘못 조합하면 가상 스프링에 저장된 탄성 에너지가
한꺼번에 풀리며 로봇이 튀어나가고, 평형점 개념을 오해하면 정상 동작을
고장으로 오독한다.
저자 역시 고가의 실물 로봇 앞에서 이런 시행착오를 겪었다.

기존 문헌이 이 어려움을 덜어 주지 못하는 이유는 두 가지다.
첫째, 유도 과정의 생략이다.
논문과 교과서는 지면 제약과 독자 가정 때문에 중간 전개를 건너뛰며, 입문자는
생략된 한 단계를 찾아 여러 문헌을 오가게 된다.
둘째, 표기의 불일치다.
같은 물리량이 문헌마다 다른 기호로 적혀, 개념이 아니라 기호 대조에 학습
시간이 소모된다.

본 총설은 이 두 문제를 정면으로 겨냥한다.
임피던스가 실제 공학 문제(19세기 장거리 전신선로의 신호 왜곡)에서 자란
언어임을 보이는 역사적 입구에서 출발하여, 라플라스 변환을 정의부터
유도하고, 전기(RLC)와 기계(질량-스프링-댐퍼) 임피던스를 하나의 일관된
표기로 연결한다.
지면 관계상 본문에는 핵심 유도를 싣고, 모든 중간 단계를 담은 확장판
원고(122쪽)와 수치 예제를 재현하는 MuJoCo 교육 시뮬레이터를 공개 저장소로
제공한다.
이어지는 2부는 이 기초 위에서 임피던스 제어, 어드미턴스 제어와의 선택
기준, 그리고 실패 사례에 기반한 안전한 파라미터 설정 원리를 다룬다.

\section{임피던스의 역사적 기원}\label{sec:u3-history}

임피던스(impedance)라는 말은 어원상 진행을 방해한다는 뉘앙스를 갖지만,
공학에서는 단순한 방해를 넘어 시스템이 입력된 흐름에 어떻게 저항하고,
에너지를 저장했다가 되돌려주는지를 함께 나타내는 양이다.
이 개념이 왜 필요했는지는 19세기 후반 장거리 전신선과 전화선의 문제로
돌아가 보면 선명해진다.

\subsection{긴 선로가 깨뜨린 저항만의 그림}

책상 위의 짧은 직류 회로에서는 흐름을 방해하는 정도를 저항 하나에 담아
옴의 법칙 $V=IR$만으로 많은 것을 설명할 수 있다.
그러나 전보는 짧은 펄스가 켜졌다 꺼지는 신호이고, 전화는 사람 목소리처럼
계속 변하는 신호라서 직류가 아니다.
실제 긴 케이블에서는 이상한 일이 벌어졌다.
멀리 보낸 펄스는 늦게 도착하며 모양이 둥글게 퍼졌고, 목소리는 높은 성분과
낮은 성분이 서로 다르게 전달되어 멀리서 뭉개져 들렸다.
펄스의 날카로운 모서리일수록 높은 주파수 성분을 많이 필요로 하므로, 선로가
높은 성분을 더 약하게 만들거나 주파수마다 도착 시간을 다르게 만들면 점과
선의 구분이 흐려진다.
전선 저항만 줄이면 된다는 직류식 생각으로는 이 현상을 설명할 수 없었다.

이유는 긴 전선이 단순한 저항이 아니기 때문이다.
긴 선로에는 단위 길이당 저항 $R'~[\Omega/\mathrm{m}]$, 인덕턴스
$L'~[\mathrm{H/m}]$, 커패시턴스 $C'~[\mathrm{F/m}]$, 절연 누설을 나타내는
컨덕턴스 $G'~[\mathrm{S/m}]$가 위치를 따라 분포한다.
$R'$와 $G'$는 에너지가 열과 누설로 사라지는 손실 요소이고, $L'$와 $C'$는
자기장과 전기장에 에너지를 잠시 저장했다가 되돌려주는 저장 요소이다.
효과를 한 점에 모아 보는 집중 소자 모델과 달리, 긴 선로는 이런 작은
효과가 위치를 따라 반복해 붙어 누적되는 분포 시스템으로 보아야 한다.
위치 $x$와 시간 $t$에서의 전압 $v(x,t)$와 전류 $i(x,t)$로 이 생각을 적으면
전신방정식(telegrapher's equations)의 기본 모양이 된다.
\begin{align}
  \frac{\partial v}{\partial x} &= -R'i - L'\frac{\partial i}{\partial t}, \\
  \frac{\partial i}{\partial x} &= -G'v - C'\frac{\partial v}{\partial t}.
\end{align}
첫째 식은 선로를 따라 갈수록 전압이 변하는 이유가 직렬 저항 손실과 직렬
인덕턴스 전압에 있다는 뜻이고, 둘째 식은 전류가 변하는 이유가 병렬 누설
전류와 커패시터 충전 전류에 있다는 뜻이다.
지금 이 식을 풀 필요는 없다.
긴 전선은 한 덩어리 저항이 아니라 작은 저장소와 손실 요소가 끝없이 이어진
물체라는 표시로만 읽으면 충분하다.

\subsection{헤비사이드와 무왜곡 조건}

올리버 헤비사이드(Oliver Heaviside)는 전보 기사로 일하며 이 문제를 직접
접했고, 맥스웰의 전자기 이론을 통신선 문제에 적용하여 긴 선로의 왜곡을
수학적으로 다루었다 \cite{heaviside_electrical_papers,donaghyspargo2018heaviside}.
그의 통찰은 신호를 가장 세게 보내는 것이 목표가 아니라, 여러 주파수 성분이
서로 비슷한 규칙으로 약해지고 지연되게 만드는 것이 목표라는 데 있다.
파형은 여러 주파수 성분의 합이므로, 일부 성분만 더 약해지거나 더 늦어지면
전체 모양이 무너진다.
균일하고 선형이며 $R'$, $L'$, $C'$, $G'$가 주파수에 거의 의존하지 않는
이상화된 선로에서, 왜곡 없는 전달 조건은
\begin{equation}
  \frac{R'}{L'} = \frac{G'}{C'}
\end{equation}
로 주어진다.
$R'/L'$는 직렬 쪽에서 손실이 자기장 저장에 비해 얼마나 큰지,
$G'/C'$는 병렬 쪽에서 누설이 전기장 저장에 비해 얼마나 큰지를 나타내며,
두 양 모두 단위가 $\mathrm{s}^{-1}$이므로 이 조건은 같은 종류의 시간 척도를
맞추는 식이다.
오래된 전화선은 누설 $G'$가 작아 $G'/C'$가 작은데 인덕턴스가 부족해
$R'/L'$가 상대적으로 컸고, 일정 간격의 코일로 유효 $L'$를 키우면 두 비율이
음성 대역에서 가까워졌다.
전류 변화를 방해하는 인덕턴스를 오히려 추가해서 신호가 좋아지는 셈인데,
이는 방해를 키우는 일이 아니라 손실과 저장의 균형을 맞추어 파형 왜곡을
줄이는 설계이다.
그래서 헤비사이드가 1880년대에 사용한 impedance라는 말은 단순한 신조어가
아니라, 저항만으로 설명되지 않는 유도성과 용량성 효과까지 함께 보며 신호
전달을 설명하려는 실무 언어였다 \cite{ethw_heaviside,ptb_impedance_history}.

\subsection{저항, 리액턴스, 임피던스}

직류에서는 흐르기 어려운 정도를 저항 하나로 말해도 충분한 경우가 많다.
그러나 교류에서는 신호가 얼마나 작아지는지뿐 아니라, 전압과 전류의
흔들림이 시간상 얼마나 앞서거나 뒤지는지도 함께 보아야 한다.
전화선 한쪽에서 한 가지 주파수의 순음을 계속 보내고 반대쪽에서 얼마나
작아졌는지, 얼마나 늦게 흔들리는지를 재는 실용적 시험을 떠올리면 된다.
인덕터는 자기장에, 커패시터는 전기장에 에너지를 저장하며 주파수에 따라
전압과 전류 사이의 위상을 바꾸는데, 이런 저장 요소가 만드는 주파수
의존적인 방해를 리액턴스(reactance)라고 부른다.
리액턴스의 단위도 옴($\Omega$)이지만, 저항처럼 평균적으로 에너지를 열로
없앤다는 뜻은 아니다.
그래서 교류의 방해는 소산 성분과 저장 성분을 함께 담아
\begin{equation}
  Z = R + \jj X
\end{equation}
로 쓴다.
실수부 $R$은 에너지를 소산하는 저항이고, 허수부 $X$가 위상을 바꾸는
리액턴스이며, $\jj$는 $90^\circ$ 위상 차이를 표시하는 허수 단위이다.
정현파 정상상태에서 세 기본 소자의 임피던스는 각각
\begin{equation}
  Z_R = R, \qquad
  Z_L = \jj\omega L, \qquad
  Z_C = \frac{1}{\jj\omega C} = -\frac{\jj}{\omega C}
\end{equation}
이다($\omega=2\pi f$는 각주파수).
인덕터는 $\omega$가 커질수록 빠른 전류 변화에 더 큰 전압을 요구하고,
커패시터는 높은 주파수 전류를 상대적으로 잘 통과시킨다.
직렬 연결에서는 세 소자에 같은 전류가 흐르고 전원이 감당할 전압은 세 소자
전압의 합이므로, 전체 임피던스도 세 임피던스의 합이 된다.
\begin{equation}
  Z(\jj\omega) = Z_R + Z_L + Z_C
  = R + \jj\left(\omega L - \frac{1}{\omega C}\right).
\end{equation}
인덕터의 양의 리액턴스와 커패시터의 음의 리액턴스는 위상 방향이 반대라서
서로 빼지며, 어떤 주파수에서는 상쇄되어 회로가 저항만 남은 것처럼 보인다.
이 상쇄 현상이 공진의 씨앗인데, 공진 자체의 자세한 전개는 공개 확장판에서 다룬다.
이렇게 하면 미분방정식으로 풀어야 할 교류 회로도 $V=IZ$라는 옴의 법칙과
닮은 대수식으로 다룰 수 있게 된다.
전신선로 해석의 세부 유도와 페이저 표기의 확장 논의는 공개 확장판을
참고한다.

전보선 이야기가 남긴 질문은 시간에 따라 변하는 입력이 시스템을 지나며
크기와 위상이 어떻게 바뀌는가이며, 전기 회로와 기계 시스템을 같은 문장으로
비교하려면 먼저 이 질문을 적을 라플라스 변환과 선형 시불변 시스템의
문법이 필요하다.


% U3 3절 — 라플라스 변환의 최소 문법 (정본: paper/sections/03_lti_system.tex 발췌 재구성)
\section{라플라스 변환의 최소 문법}
\label{sec:u3-laplace}

앞 절의 전신선로가 남긴 질문 --- 신호의 크기와 위상이 왜, 어떻게 변하는가 ---
을 계산 가능한 형태로 바꾸려면 두 가지 준비가 필요하다.
하나는 시스템에 대한 가정이고, 다른 하나는 미분방정식을 대수식으로 바꾸는
변환이다.
이 절은 그 최소한의 문법을 유도 생략 없이 정리한다.

먼저 가정부터 보자.
시스템이 선형(linear)이라는 것은, 입력 $u(t)$를 출력 $y(t)$로 바꾸는 규칙
$\mathcal{S}$가 임의의 입력 $u_1$, $u_2$와 스칼라 $a$, $b$에 대해
$\mathcal{S}\{a u_1 + b u_2\} = a\,\mathcal{S}\{u_1\} + b\,\mathcal{S}\{u_2\}$를
만족한다는 뜻이다.
입력을 두 배로 하면 출력도 두 배가 되고, 두 입력을 더해 넣은 결과는 각각의
출력을 더한 것과 같다.
시불변(time-invariant)이라는 것은 같은 입력을 $t_0$만큼 늦게 넣으면 출력도
똑같이 $t_0$만큼 늦어질 뿐 모양은 변하지 않는다는 뜻,
즉 $\mathcal{S}\{u(t-t_0)\} = y(t-t_0)$이다.
실제 로봇은 완전한 선형도 시불변도 아니지만, 작은 움직임과 한 작동점
근처에서는 이 선형시불변(LTI) 가정이 많은 현상을 충분히 잘 설명하며,
전달함수와 주파수 응답이라는 강력한 도구를 쓸 수 있게 해 준다.

\subsection{정의에서 미분 규칙까지}

라플라스 변환은 시간 함수 $x(t)$를 복소 변수 $s$의 함수로 요약하는 적분이다.
\begin{equation}
  X(s) = \mathcal{L}\{x(t)\}
  = \int_{0}^{\infty} x(t)\,e^{-st}\,\dd t.
  \label{eq:u3-laplace-def}
\end{equation}
말로 풀면, $x(t)$에 $e^{-st}$라는 점점 작아지는 추를 곱해 $t=0$부터 끝까지
모두 더한 값이다.
$s$를 하나 고를 때마다 숫자가 하나 나오므로, 시간 파형 전체가 $X(s)$ 한
장으로 요약된다.

이 변환이 유용한 이유는 미분이 $s$를 곱하는 연산으로 바뀌기 때문인데,
이 규칙은 외울 공식이 아니라 정의 \eqref{eq:u3-laplace-def}와 곱의 미분
규칙만으로 나온다.
$s$는 $t$와 무관하므로
\begin{equation}
  \frac{\dd}{\dd t}\left[x(t)\,e^{-st}\right]
  = \dot{x}(t)\,e^{-st} - s\,x(t)\,e^{-st}
\end{equation}
이다.
양변을 $t=0$부터 $\infty$까지 적분하자.
왼쪽은 미분한 것을 다시 적분하므로 끝값 빼기 시작값,
즉 $x(t)e^{-st}$의 $t\to\infty$ 값에서 $t=0$ 값 $x(0)$을 뺀 것이다.
응답이 발산하지 않아 $t\to\infty$에서 $x(t)e^{-st}\to 0$이 되는 $s$에서 보면
왼쪽은 $0 - x(0) = -x(0)$이다.
오른쪽의 두 적분은 정의에 따라 각각 $\mathcal{L}\{\dot{x}\}$와 $-sX(s)$이므로
$-x(0) = \mathcal{L}\{\dot{x}\} - sX(s)$, 정리하면
\begin{equation}
  \mathcal{L}\{\dot{x}(t)\} = sX(s) - x(0)
  \label{eq:u3-deriv-rule}
\end{equation}
이다.
초기조건을 0으로 두면 $\dot{x}$는 $sX(s)$처럼 다룰 수 있다.

간단한 검산을 해 보자.
$x(t)=1$(상수)이면 정의에서
$X(s)=\int_0^\infty e^{-st}\dd t
=\left[-e^{-st}/s\right]_0^\infty
=0-\left(-\tfrac{1}{s}\right)=\tfrac{1}{s}$이다.
$\dot{x}=0$이므로 규칙 \eqref{eq:u3-deriv-rule}의 왼쪽은
$\mathcal{L}\{0\}=0$이고, 오른쪽도
$sX(s)-x(0)=s\cdot(1/s)-1=0$으로 일치한다.

두 번 미분한 항은 같은 규칙을 두 번 적용하면 된다.
$\dot{x}$를 새 함수로 보고 \eqref{eq:u3-deriv-rule}를 다시 쓰면
\begin{equation}
  \mathcal{L}\{\ddot{x}\}
  = s\,\mathcal{L}\{\dot{x}\} - \dot{x}(0)
  = s^2 X(s) - s\,x(0) - \dot{x}(0)
\end{equation}
이고, 초기조건을 0으로 두면 $\ddot{x}$는 $s^2 X(s)$처럼 다룰 수 있다.

\subsection{미분방정식에서 전달함수로}

이제 질량-스프링-댐퍼의 운동방정식 $m\ddot{x}+b\dot{x}+kx=f$의 각 항을
초기조건 0에서 하나씩 바꾸면
$m\ddot{x}\to m s^2 X(s)$, $b\dot{x}\to b s X(s)$, $kx\to kX(s)$,
$f\to F(s)$이므로
\begin{equation}
  (m s^2 + b s + k)\,X(s) = F(s)
\end{equation}
가 된다.
시간에 따라 계속 갱신해야 하는 미분 장부가 한 장의 대수 장부로 바뀐 것이다.
입력과 출력의 비율로 정리하면 전달함수(transfer function)
\begin{equation}
  H(s) = \frac{X(s)}{F(s)} = \frac{1}{m s^2 + b s + k}
  \label{eq:u3-msd-tf}
\end{equation}
를 얻는다.
분모의 $ms^2$ 항은 관성, $bs$ 항은 감쇠, $k$ 항은 스프링 복원의 영향을 담는다.
전달함수는 초기 저장 에너지를 0으로 둔(zero-state) 입력-출력 관계이며,
같은 물리 장치라도 무엇을 입력과 출력으로 고르느냐에 따라 다른 전달함수로
보인다는 점만 기억해 두자.

여기서 $s$는 단순한 알파벳이 아니라 시간응답의 성질을 담는 복소 변수로,
보통 $s=\sigma+\jj\omega$로 생각한다.
$\sigma$는 응답이 커지거나 줄어드는 속도, $\omega$는 흔들리는 속도와
관련된다.
실제 물리 입력이 복소수라는 뜻이 아니라, 감쇠와 진동을 한 기호 안에서 함께
계산하기 위한 표현이다.

\subsection{주파수 응답 읽기}

안정한 LTI 시스템에 정현파 입력 $u(t)=A\cos(\omega t)$를 넣으면, 정상상태
출력은 같은 주파수를 가지며 크기와 위상만 달라진다.
\begin{equation}
  y_{\mathrm{ss}}(t)
  = A\,|H(\jj\omega)|\cos\!\left(\omega t + \angle H(\jj\omega)\right).
\end{equation}
$|H(\jj\omega)|$와 $\angle H(\jj\omega)$는 복소수 $H(\jj\omega)$의 크기와
각도이다.
복소수 $a+\jj b$를 평면 위의 점 $(a,b)$로 그리면 크기는 원점까지의 거리
$\sqrt{a^2+b^2}$, 각도는 양의 실수축에서 잰 회전각이다.
예를 들어 어떤 주파수에서 $H(\jj\omega)=1-\jj$라면 크기는
$\sqrt{1^2+(-1)^2}=\sqrt{2}$, 각도는 $-45^\circ$이므로, 출력 진폭은 입력의
약 1.41배가 되고 위상은 $45^\circ$만큼 늦는다.
즉 $s=\jj\omega$를 대입한 $H(\jj\omega)$ 하나로, 각 주파수에서 시스템이
얼마나 증폭하거나 감쇠하는지, 얼마나 늦게 반응하는지를 모두 읽을 수 있다.
pole과 응답 모양의 관계, 컨볼루션, 최종값 정리 등 여기서 줄인 내용의 세부
유도는 확장판을 참고한다.

이제 이 문법을 첫 번째 실물에 적용할 차례다.
다음 절에서는 저항, 인덕터, 커패시터로 이루어진 RLC 회로의 방정식을 같은
절차로 변환하여, 전기 임피던스가 어떻게 읽히는지 본다.


% U3 4절 — 전기 임피던스: RLC 회로 (정본: paper/sections/04_electric_system.tex 발췌 재구성)
\section{전기 임피던스: RLC 회로}
\label{sec:u3-electric}

앞 절에서 준비한 문법을 첫 번째 실물에 적용한다.
저항 $R$, 인덕터 $L$, 커패시터 $C$가 한 전류 경로에 놓인 직렬 RLC 회로는
2차 미분방정식으로 쓸 수 있는 가장 작은 전기 시스템이다.
세 소자의 역할에서 시작해 3절의 규칙으로 전달함수와 표준 2차 형식까지
유도 생략 없이 내려간다.

\subsection{세 소자의 역할: 소산과 두 저장}

저항은 전류가 지나가는 동안 전기 에너지를 열로 바꾸는 소자이다.
전류가 들어가는 단자를 전압의 $+$로 잡는 수동부호약속에서 $v_R = R i$이므로,
저항이 흡수하는 순간 전력은 $p_R = v_R\,i = R\,i^2$이다.
전류 방향이 바뀌어도 $i^2$는 항상 양수이므로, 저항은 에너지를 되돌려주지
않고 계속 소산한다.

인덕터는 전류가 만드는 자기장에 에너지를 저장하는 소자로,
\begin{equation}
  v_L(t) = L\frac{\dd i(t)}{\dd t}
  \label{eq:u3-ind-v}
\end{equation}
를 만족한다.
전류를 빠르게 바꿀수록 큰 전압이 필요하므로, 인덕터는 전류의 변화에 버틴다.
순간 전력에 \eqref{eq:u3-ind-v}를 넣고 곱의 미분 규칙으로 정리하면
$p_L = v_L\,i = L\,i\frac{\dd i}{\dd t}
= \frac{\dd}{\dd t}\bigl(\tfrac{1}{2}L i^2\bigr)$이므로,
인덕터에 들어간 전력은 열로 사라지는 대신 자기장 에너지
$E_L = \frac{1}{2}Li^2$의 변화율로 쌓인다.
전류가 줄어드는 구간에서는 $p_L<0$이 되어 저장했던 에너지를 회로로 되돌려준다.

커패시터는 반대 방향의 저장 소자로, 분리된 전하 $q = C v_C$에서 전류는
전하의 변화율이므로
\begin{equation}
  i_C(t) = C\frac{\dd v_C(t)}{\dd t}
  \label{eq:u3-cap-i}
\end{equation}
이고, 전압을 빠르게 바꾸려면 많은 전하를 짧은 시간에 옮겨야 하므로 커패시터는 전압의 변화에 버틴다.
같은 곱의 미분 계산을 반복하면
$p_C = v_C\,i_C = C\,v_C\frac{\dd v_C}{\dd t}
= \frac{\dd}{\dd t}\bigl(\tfrac{1}{2}C v_C^2\bigr)$이므로,
커패시터의 저장 에너지는 전기장 에너지 $E_C = \frac{1}{2}Cv_C^2$이다.
정현파에서 미분이 $\jj\omega$ 곱으로 바뀐다는 3절의 관찰을 적용하면
세 소자의 임피던스는 $Z_R = R$, $Z_L = \jj\omega L$, $Z_C = 1/(\jj\omega C)$이고,
$\jj$의 유무가 곧 소산과 저장의 차이, 전압-전류 위상이 $\pm 90^\circ$
어긋나는가의 차이로 나타난다.

\subsection{직렬 RLC 방정식과 전달함수}

회로를 한 방향으로 도는 기준 전류를 $i(t)$, 그 전류가 들어가는 커패시터
판의 전하를 $q(t)$라 하면 $i = \dot{q}$이다.
키르히호프 전압 법칙에 따라 한 바퀴를 돌며 만나는 전압 변화의 합은 0,
즉 $v_{\mathrm{in}} - v_L - v_R - v_C = 0$이다.
각 소자 관계 $v_L = L\ddot{q}$, $v_R = R\dot{q}$, $v_C = q/C$를 대입하면
\begin{equation}
  L\ddot{q}(t) + R\dot{q}(t) + \frac{1}{C}q(t) = v_{\mathrm{in}}(t)
  \label{eq:u3-rlc-ode}
\end{equation}
를 얻는다.
이 한 줄에서 $L\ddot{q}$는 전류 변화에 대한 인덕터 전압, $R\dot{q}$는
소산 전압, $q/C$는 쌓인 전하에 비례해 커지는 복원 성격의 전압이다.

이제 3절의 미분 규칙 \eqref{eq:u3-deriv-rule}대로 영 초기조건에서 각 항을
하나씩 바꾼다.
$L\ddot{q} \to L s^2 Q(s)$, $R\dot{q} \to R s Q(s)$,
$q/C \to Q(s)/C$, $v_{\mathrm{in}} \to V_{\mathrm{in}}(s)$이므로
\begin{equation}
  \left(L s^2 + R s + \frac{1}{C}\right) Q(s) = V_{\mathrm{in}}(s),
  \qquad
  \frac{Q(s)}{V_{\mathrm{in}}(s)} = \frac{1}{L s^2 + R s + \frac{1}{C}}
  \label{eq:u3-rlc-tf}
\end{equation}
이다.
전하를 출력으로 잡으면 $i = \dot{q}$와 $v_C = q/C$로 전류와 커패시터
전압을 바로 얻을 수 있고, 출력 선택이 바뀌어도 분모는 같다.

\subsection{표준 2차 형식과 두 파라미터}

서로 다른 물리 시스템을 같은 눈금으로 비교하기 위해 2차 시스템은 보통
\begin{equation}
  \ddot{x} + 2\zeta\omega_n\dot{x} + \omega_n^2 x = \omega_n^2 u
  \label{eq:u3-std-2nd}
\end{equation}
로 쓴다.
$\omega_n$은 고유진동수, $\zeta$는 감쇠비이고, $x$는 두 번 미분되는 대표
상태를 임시로 부르는 이름이다.
식 \eqref{eq:u3-rlc-ode}의 양변을 $L$로 나누면
\begin{equation}
  \ddot{q} + \frac{R}{L}\dot{q} + \frac{1}{LC}q = \frac{1}{L}v_{\mathrm{in}}
\end{equation}
이고, \eqref{eq:u3-std-2nd}와 항별로 계수를 비교하면
\begin{equation}
  2\zeta\omega_n = \frac{R}{L},
  \qquad
  \omega_n^2 = \frac{1}{LC},
  \qquad
  \omega_n^2 u = \frac{1}{L}v_{\mathrm{in}}
\end{equation}
이다.
마지막 식에 $\omega_n^2 = 1/(LC)$를 넣으면 $u = C v_{\mathrm{in}}$,
즉 표준형의 입력은 전압에 커패시턴스를 곱한 전하 단위의 계산용 입력으로,
표준형이 물리가 아니라 눈금을 바꾸는 일임이 여기서 잘 보인다.
위치항에서 $\omega_n = 1/\sqrt{LC}$이고, 속도항 비교에서 $\zeta$만 남기면
\begin{equation}
  \zeta
  = \frac{R/L}{2\omega_n}
  = \frac{R}{2L}\sqrt{LC}
  = \frac{R}{2}\sqrt{\frac{LC}{L^2}}
  = \frac{R}{2}\sqrt{\frac{C}{L}}
\end{equation}
이다.
저항을 키우면 감쇠비가 커지고, $L$과 $C$는 자기장 에너지와 전기장
에너지가 서로 오가는 리듬 $\omega_n$을 정한다.

숫자로 확인해 보자.
$L = 10~\mathrm{mH} = 0.01~\mathrm{H}$,
$C = 100~\mu\mathrm{F} = 0.0001~\mathrm{F}$이고, 임계감쇠 $\zeta = 1$에
해당하는 저항을 $R_c$라 하면
\begin{equation}
  \omega_n = \frac{1}{\sqrt{LC}} = 1000~\mathrm{rad/s},
  \quad
  f_n = \frac{\omega_n}{2\pi} \approx 159~\mathrm{Hz},
  \quad
  R_c = 2\sqrt{\frac{L}{C}} = 20~\Omega
\end{equation}
이다.
$R$이 이보다 작으면 부족감쇠로 흔들리며 수렴하고, 크면 과감쇠 쪽으로 간다.
LC 자유진동의 에너지 왕복, 직렬 공진 $\omega L = 1/(\omega C)$와 튜닝
응용, 소자들의 역사적 배경 등 여기서 줄인 내용은 공개 확장판을 참고한다.

다음 절에서는 같은 2차 구조가 눈에 보이는 운동으로 다시 나타나,
$L\ddot{q} + R\dot{q} + q/C = v_{\mathrm{in}}$의 각 자리에
$m\ddot{x} + b\dot{x} + kx = f$가 그대로 겹쳐지는 것을 본다.


% U3 5절 — 기계 임피던스: 질량-스프링-댐퍼 (정본: paper/sections/05_mechanical_system.tex 발췌 재구성)
\section{기계 임피던스: 질량-스프링-댐퍼}
\label{sec:u3-mechanical}

앞 절의 문법을 첫 실물에 적용했으니, 이제 같은 2차 구조를 눈에 보이는 운동으로 다시
읽는다.
질량-스프링-댐퍼(mass-spring-damper)는 진동과 임피던스를 설명하는 가장 작은 기계
모델이다.
질량 $m$은 속도 변화를 싫어하는 관성, 강성 $k$는 변위를 평형점으로 되돌리는 복원력의
기울기, 감쇠계수 $b$는 속도에 비례해 운동 에너지를 빼내는 저항이다.
전기 쪽으로 보면 질량은 인덕터, 스프링은 커패시터, 댐퍼는 저항에 대응한다.

\subsection{운동방정식: 힘의 장부}

평형점을 $x=0$으로 두고 외력을 $f(t)$라 하자.
물체에 실제로 작용하는 스프링 힘과 댐퍼 힘은 변위와 속도를 줄이는 방향이므로
$f_s=-kx$, $f_d=-b\dot{x}$이다.
뉴턴의 제2법칙 \cite{newton1687principia}을 $m\ddot{x}=f+f_s+f_d=f-kx-b\dot{x}$로 쓰고(스프링 힘 $f_s=-kx$는 훅의 법칙 \cite{hooke1678depotentiarestitutiva}) 스프링항과 감쇠항을
왼쪽으로 옮기면
\begin{equation}
  m\ddot{x}(t)+b\dot{x}(t)+kx(t)=f(t)
  \label{eq:u3m-msd}
\end{equation}
가 된다.
이는 외울 공식이라기보다 외력 $f$가 감당할 부담을 관성항 $m\ddot{x}$, 감쇠항
$b\dot{x}$, 스프링항 $kx$로 나누어 적은 장부이다.

\subsection{고유진동수}

감쇠와 외력이 없는 $m\ddot{x}+kx=0$, 즉 $\ddot{x}+(k/m)x=0$을 먼저 본다.
감쇠 없는 왕복 운동은 일정한 모양으로 반복되므로 $x(t)=A\cos(\omega t)$로 시험한다.
두 번 미분하면 $\ddot{x}=-\omega^2 x$이므로 $-m\omega^2 x+kx=0$, $x\neq0$인 진동에서는
$\omega^2=k/m$, 따라서 $\omega_n=\sqrt{k/m}$이다.
단위로도 확인된다.
\begin{equation}
  \frac{\mathrm{N/m}}{\mathrm{kg}}
  =
  \frac{\mathrm{kg\,m/s^2}/\mathrm{m}}{\mathrm{kg}}
  =
  \frac{1}{\mathrm{s^2}}
\end{equation}
이므로 제곱근을 취하면 $1/\mathrm{s}$가 되고, rad를 무차원으로 취급해 rad/s라 부른다.
고유진동수는 되돌리려는 성질과 둔감함의 비율이라, 강성 4배는 $\omega_n$ 2배, 질량 4배는
$\omega_n$ 절반이라는 제곱근 관계를 가진다.

\subsection{표준형과 계수 맞추기}

식 \eqref{eq:u3m-msd}의 양변을 $m$으로 나누고 표준 2차 시스템
\begin{equation}
  \ddot{x}(t)+2\zeta\omega_n\dot{x}(t)+\omega_n^2 x(t)=\frac{1}{m}f(t)
  \label{eq:u3m-standard}
\end{equation}
과 계수를 견주면, $x$ 항에서 $\omega_n^2=k/m$, $\dot{x}$ 항에서 $2\zeta\omega_n=b/m$이므로
\begin{equation}
  \omega_n=\sqrt{\frac{k}{m}}, \qquad \zeta=\frac{b}{2\sqrt{mk}}
\end{equation}
이다.
$\omega_n^2$은 위치를 되돌리는 항의 세기, $2\zeta\omega_n$은 속도에 비례해 진동을 줄이는
항의 세기이다.
감쇠비 $\zeta$가 무차원인 것은 뒤의 임계감쇠 $b_c=2\sqrt{mk}$로 $b$를 나눈 값이라,
같은 댐퍼라도 $m$과 $k$가 다르면 효과가 달라진다는 사실을 한 숫자에 담기 때문이다.

\subsection{특성근과 응답 형태}

자유응답 $f=0$에서 $x(t)=e^{st}$를 넣으면 특성방정식 $ms^2+bs+k=0$을 얻고,
$s_{1,2}=(-b\pm\sqrt{b^2-4mk})/(2m)$이다.
논문에서 한 줄로 건너뛰는 표준 파라미터 대입을 두 조각으로 나눠 보자.
$\zeta=b/(2\sqrt{mk})$를 뒤집으면 $b=2\zeta\sqrt{mk}$이므로 분수 앞부분은
\begin{equation}
  -\frac{b}{2m}=-\frac{2\zeta\sqrt{mk}}{2m}=-\zeta\sqrt{\frac{k}{m}}=-\zeta\omega_n,
\end{equation}
근호 안은 $b^2=4\zeta^2 mk$이므로
\begin{equation}
  b^2-4mk=4\zeta^2 mk-4mk=4mk(\zeta^2-1)
\end{equation}
이다.
$\sqrt{b^2-4mk}=2\sqrt{mk}\sqrt{\zeta^2-1}$을 $2m$으로 나누면 $\omega_n\sqrt{\zeta^2-1}$
이 되어, 두 조각을 합치면
\begin{equation}
  s_{1,2}=-\zeta\omega_n\pm\omega_n\sqrt{\zeta^2-1}
  \label{eq:u3m-roots}
\end{equation}
이다.
$m=1$, $b=2$, $k=4$로 검산하면 $\omega_n=\sqrt{4}=2$, $\zeta=2/(2\sqrt{4})=0.5$이고,
공식으로는 $s_{1,2}=(-2\pm\sqrt{4-16})/2=-1\pm\sqrt{-3}$, 표준형으로도 $-\zeta\omega_n=-1$,
$\omega_n\sqrt{\zeta^2-1}=2\sqrt{-0.75}=\sqrt{-3}$이라 정확히 일치한다.

부족감쇠 $0<\zeta<1$에서는 근호 안이 음수가 되어 근이 복소수이다.
$\jj$가 제곱하면 $-1$인 허수 단위이므로 $\zeta^2-1=(-1)(1-\zeta^2)$로 쪼개면
$\sqrt{\zeta^2-1}=\jj\sqrt{1-\zeta^2}$이고($1-\zeta^2>0$이라 마지막은 보통의 양수 제곱근),
특성근은 $s_{1,2}=-\zeta\omega_n\pm\jj\omega_d$가 된다.
여기서 $\omega_d=\omega_n\sqrt{1-\zeta^2}$는 감쇠 진동수(damped natural frequency)이며,
부족감쇠 자유응답은 줄어드는 포락선 $e^{-\zeta\omega_n t}$ 안에서 $\omega_d$로 흔들리는
$x(t)=e^{-\zeta\omega_n t}(A\cos\omega_d t+B\sin\omega_d t)$ 형태가 된다.

\subsection{감쇠 영역과 임계감쇠}

실수부 $-\zeta\omega_n$은 응답이 공통으로 줄어드는 속도를, 제곱근 항은 근이 실수인지
복소수인지를 가른다.
판별식 $b^2-4mk$가 음수이면 흔들리고 양수이면 흔들리지 않으므로, 경계는
$b^2-4mk=0$, 즉 $b=\pm2\sqrt{mk}$이다.
감쇠계수는 물리적으로 양수이므로 양수 근을 택해 $b_c=2\sqrt{mk}$가 된다.
$\zeta=0$이면 진동이 계속되고, $0<\zeta<1$이면 흔들리며 줄어들며, $\zeta=1$이면 두 근이
$s=-\omega_n$으로 겹쳐 반복 진동 없이 돌아오는 경계 응답이 되고, $\zeta>1$이면 서로 다른
두 음의 실수근 중 더 느린 항이 전체 수렴을 지배해 오히려 둔해질 수 있다.

\subsection{오버슈트와 정착시간}

일정한 힘 $F_0$가 계속 걸리면 오래 뒤 $\dot{x}=\ddot{x}=0$이 되어 스프링 항만 남으므로
최종값은 $x_{\mathrm{ss}}=F_0/k$이다.
이 값을 얼마나 지나치는지가 오버슈트이며, 표준 2차 시스템, 영 초기조건, 양의 단위계단
입력, $0<\zeta<1$ 조건에서 최대 오버슈트 비율은
\begin{equation}
  M_p=\exp\!\left(-\frac{\zeta\pi}{\sqrt{1-\zeta^2}}\right)
  \label{eq:u3m-overshoot}
\end{equation}
이다.
$\zeta$가 작으면 남은 에너지가 많아 크게 지나치는데, $\zeta=0.2$이면 약 $53\%$를
예측한다.
$\zeta=0.7$이면 $\sqrt{1-\zeta^2}=\sqrt{0.51}\approx0.714$, 지수 안은
$-0.7\pi/0.714\approx-3.08$이므로 $M_p=e^{-3.08}\approx0.046$, 즉 약 $4.6\%$로 줄어든다.

정착시간은 응답이 목표 주변의 작은 띠 안에 들어간 뒤 계속 머무는 시간이다.
진폭 포락선이 대략 $e^{-\zeta\omega_n t}$처럼 줄어든다고 보고 $\pm2\%$ 기준, 즉
$0.02=1/50$까지 줄어드는 시간을 잡으면 $e^{-\zeta\omega_n t}\approx0.02$, 양변에 로그를
취해 $\zeta\omega_n t\approx\ln 50\approx3.9$가 된다.
$e^{3.9}\approx50$이므로 포락선이 약 50분의 1로 줄었다고 읽을 수 있고, 여기서
$T_s\approx 4/(\zeta\omega_n)$라는 기억하기 쉬운 근사가 나온다.
$\omega_n=10$, $\zeta=0.7$이면 $T_s\approx4/(0.7\cdot10)\approx0.57~\mathrm{s}$이다.
이 근사는 대략 $0.4\lesssim\zeta\lesssim0.8$에서 유용하며, 5\% 기준을 쓰면 상수가 대략
3에 가까워진다.
힘의 관점으로 본 진동을 운동 에너지와 탄성 퍼텐셜 에너지의 교환으로 다시 읽는 에너지
장부와 응답 그래프의 세부는 확장판을 참고한다.

\subsection{기계적 임피던스}

같은 운동방정식이라도 무엇을 출력과 비율로 잡느냐에 따라 이름이 달라진다.
식 \eqref{eq:u3m-msd}를 영 초기조건에서 라플라스 변환하면 $F(s)=(ms^2+bs+k)X(s)$,
속도는 $V_x(s)=sX(s)$이므로 힘을 속도로 나누면 기계적 임피던스
\begin{equation}
  Z_m(s)=\frac{F(s)}{sX(s)}=ms+b+\frac{k}{s}
  \label{eq:u3m-impedance}
\end{equation}
를 얻는다.
$s=\jj\omega$를 넣으면 $Z_m(\jj\omega)=b+\jj(\omega m-k/\omega)$로, 감쇠항 $b$는 주파수와
무관한 소산 성분으로 남고, 질량항 $\omega m$은 빠르게 흔들수록, 스프링항 $k/\omega$는
천천히 밀어 위치를 바꿀수록 커진다.
2절에서 전기 임피던스 $Z=R+\jj(\omega L-1/\omega C)$를 예고편으로 두었던 자리에, 이제
저항 대신 댐퍼, 인덕턴스 대신 질량, 커패시턴스 대신 스프링이 같은 형태로 들어선다.
임피던스는 얼마나 딱딱한가만이 아니라, 주파수에 따라 스프링처럼, 질량처럼, 감쇠처럼
다르게 보이는 힘-운동 관계임을 보여준다.

이 2차 언어 --- $\omega_n$, $\zeta$, 그리고 강성ㆍ감쇠ㆍ관성으로 나뉘는 임피던스 --- 가
바로 다음 절에서 로봇의 접촉 거동을 설계할 때 쓰는 그 언어이다.


% U3 6~7절 — 로봇 접촉 제어 전망 + 결론
\section{로봇 접촉 제어를 향해}
\label{sec:u3-outlook}

지금까지의 여정을 한 줄로 줄이면 이렇다.
임피던스는 힘과 흐름 사이의 동역학적 관계이고, 그 관계는 전기에서도
기계에서도 같은 2차 구조 --- 저장 두 개와 소산 하나 --- 로 나타난다.
로봇 접촉 제어는 이 언어의 자연스러운 다음 장면이다.
로봇 손끝에 가상의 스프링과 댐퍼를 달아 원하는 힘-운동 관계를 설계하는
임피던스 제어에서, 강성 \(k_d\)와 감쇠비 \(\zeta\)는 본 총설에서 준비한
바로 그 눈금으로 읽힌다 \cite{hogan1985impedancepart1}.
예를 들어 접촉에서 허용할 힘과 밀림을 정하면 강성의 출발값이 나오고,
목표점을 멀리 잡은 채 강성을 켜는 순간 \(\frac{1}{2}k_d\delta^2\)의
에너지가 이미 장전된다는 계산은 5절의 스프링 에너지 감각 그대로이다.

이 다음 장면 --- 임피던스 제어의 유도, 어드미턴스 제어와의 선택 기준,
그리고 실제 사고 사례에 기반한 안전한 파라미터 설정 원리 --- 는 본 총설의
2부에서 다룬다.
본 1부와 2부의 모든 수치 예제는 공개 MuJoCo \cite{todorov2012mujoco} 교육
시뮬레이터에서 재현할 수 있다.
예를 들어 저장 에너지를 계산하지 않고 큰 강성을 명령했을 때의 발사 사고는
저장소의 \texttt{configs/lab01\_msd/f2\_launch\_high\_energy.yaml} 설정
하나로 안전하게 겪어볼 수 있다.
전체 유도를 담은 확장판 원고, 시뮬레이터 코드, 실행 가이드는 공개
저장소\footnote{\url{https://github.com/ycpiglet/manipulator-control-tutorial}}에서
제공한다.

\section{결론}
\label{sec:u3-conclusion}

본 총설 1부는 임피던스라는 언어를 유도 생략 없이, 하나의 표기로, 역사적
동기와 함께 정리했다.
장거리 전신선로의 신호 왜곡이라는 실제 문제에서 출발해, 라플라스 변환을
정의에서부터 유도하고, RLC 회로와 질량-스프링-댐퍼가 같은 2차 구조를
공유함을 계수 비교와 수치 검산으로 확인했다.
독자가 이 여정에서 얻어야 할 것은 공식 목록이 아니라, 고유진동수와
감쇠비라는 두 눈금으로 어떤 2차 시스템이든 읽어내는 감각이다.
2부는 이 감각을 로봇 손끝으로 옮겨, 접촉 제어 방식의 선택과 안전한
파라미터 설정을 다룬다 \cite{lynchpark2017modernrobotics}.


\bibliographystyle{unsrt}
\bibliography{../../references/refs}
\end{document}
